Vele sterren flonkerden
aan de heldere hemel.
Mistwolken verdonkerden.
Een rover schoot een kemel

toen ze vergeefs haar schaduw
poogde te integreren.
Maar haar zelfkant bleek te sluw.
Tot ze kon accepteren

zijn ware betekenis.
Zijn taak was te verbergen
haar onkiese duisternis.
Het zou veel begrip vergen

en het doorstaan van groeipijn
om dit kerndeel te hechten
aan haar begrensd bewustzijn
en de schaduw te knechten.

De speelse avonturier
naaide net als Peter Pan
op een tedere manier
de aangenomen Satan

aan haar zelf, permanent vast.
Zo werd de persona heel
en de ongewenste last
sindsdien een kostbaar juweel.

Het sieraad keek naar een boom
met twee paradijsvogels.
De een was telkens op stoom,
prikkels vurend als kogels

en vruchten etend naar gril,
rusteloos met veel gesmak.
De ander zat vredig stil
op een hogere boomtak,

die haar een weids uitzicht bood.
Hun veren glansden een tint,
levendig, vurig en rood,
als een kampvuur in de wind.

Kleine donkere stippen
bedekten hun borst en kop.
De laagste bleef maar wippen
van tak naar tak als schrokop.

Een paar vruchten waren zoet,
anderen bleken bitter.
Soms voelde hij zich dus goed,
maar vaak wat minder fitter.

Kijkend door deze ogen,
vroeg de vrijbuiter zich af
hoe zijn peer onbewogen
zich aan alles overgaf.

Zonder maar iets te eten,
toch altijd tevreden leek.
De laagste moest dit weten!
Tot zij naar de hoogste keek

en zag dat ze eender zijn;
de universele ziel.
Getroffen door lust en pijn
in Jiva’s achilleshiel.

Zo was haar temperament,
in voldoende bewustzijn
en indien op tijd herkend,
wellicht een bron van welzijn,

om emotionaliteit
snel te laten afvlakken.
Haar eigen intensiteit
kon haar niet meer inpakken.

"Mama, waar blijf je?"; zei ze
op een matras van veren.
Men noemde haar Agape
om haar goedheid te eren.

Altruïstische liefde,
die ze goedgelovig gaf,
bombardeerde Agape,
volgens de Griekse maatstaf,

tot het hoogste ideaal
wat betreft menslievendheid.
Haar liefde was niet zo schraal
als familiariteit,

verliefdheid of een vriendschap.
Ook al was ze nog zo jong,
ze toonde haar gemeenschap
de werkelijke oorsprong

voor broederschap en eenheid.
Niemand werd geslachtofferd.
Een besef van gelijkheid,
waardoor men zich opoffert.

Ze deelde haar liefdesmaal
met behoeftige armen.
Ieders welzijn stond centraal.
De basis voor erbarmen.

Ze was van haar ego los,
maar zeker niet egoloos.
Eén zelf als de kosmos.
Ze was soms blij en soms boos.

De basis voor echt gezag.
Zo verwierf ze alom faam
met haar heilige gedrag.
Liebling werd zo haar koosnaam.

Ze negeerde haar status,
ook al droomde Agape
in een luxueus en knus
hemelbed onder zijde.

Gordijnen hingen rondom,
die privacy ophouden.
Ze waren ook geschikt om
haar warmte vast te houden.

Een rode haar in een kam,
liggend op de kaptafel,
werd verlicht door een kaarsvlam,
die stond op de schrijftafel.

Daar werd ook een schets verlicht,
typisch uit kinderjaren,
met een moederlijk gezicht,

omringd door rode haren.
Ze bezat weinig kleding,
ondanks de ruime kasten.
Ze droeg ze met waardering

tot ze niet meer goed pasten
en ze een arm kind het schonk.
Zo ook haar berg juwelen,
die na de gift snel weer slonk
als iemand kwam bedelen

en zij hem iets kon schenken.
Zoals op haar verjaardag
- met veel luxe geschenken -
die net achter Liebling lag.

Een dag met haar familie
vol met festiviteiten.
Toch toonde ze rebellie
door vooral tijd te slijten

met veelminderbedeelden.
Toen een groep muzikanten
tijdens het jaarfeest speelden,
ontving Liebling passanten.

Ze gaf delicatessen
uit de grootkeuken gepakt
en onthaalde hun lessen. Mensen, die leken geknakt

te zijn door een zwaar leven,
bleken veel levensvreugde
en geestkracht te beleven.
Wat de kleine meid heugde,

omdat haar gangbaar bezoek,
meestal geen passie meebracht.
Die geestdrift was tevens zoek
als ze de kerkdienst doorbracht.

Agapa werd vooral moe
van al het geestdodende,
zielloos en ijdel gedoe.
Dat was het verstorende

van wat haar energie gaf.
Dus volgde al snel een gaap.
Zo viel Agape bekaf
in een welverdiende slaap.

Liebling’s gouvernante
waakte in een zitmeubel
scherp over de charmante.
Ze schoof met wat gefröbel,

na Agape’s slaappraten,
haar dikke deken terug.
Ze kon het niet nalaten
en gaf een kus op haar rug.

Op de vloer lag dik tapijt,
die koude voetjes strelen.
Het stucplafond was gewijd
aan geloofstaferelen.

Schilderijen die tevens
alle wanden versierden.
Afwisselend stillevens
met portretten, die vierden

haar nietige voorouders.
Het gros was het onwaardig.
Geen bezat brede schouders.
Nooit was er eentje aardig.

“Geeft niet! Volgend jaar beter.”
De gouvernante aaide
de dromende debater,
die zich daarop omdraaide.

In dromen zag ze beelden
van tijden lang vervlogen.
Herinneringen speelden
en gevoelens bewogen.

Ze zag haar vader knielen
naast haar in een grasweide,
met ogen die bevielen
en een lach die verblijdde.

De gouvernante naast hem,
had ook trouw in haar ogen.
Liebling hoorde in haar stem
liefde en mededogen.

Ze renden door het grasveld,
geparfumeerd door mestgeur.
In dromen was haar wereld
een plek zonder pijn of sleur.

Liebling rolde in 't gras,
met veldbloemen om haar heen.
Haar vader volgde al ras
en knielde bij haar sereen.

Hij wilde ervoor zorgen
dat Agape hem niet vreest.
Ze voelde zich geborgen.
Hij begreep haar kindergeest.

Ook nadat ze had gehuild,
omdat ze door haar gerol
haar nieuwe jurk had bevuild
met een flinke koeiendrol.

Hij erkende haar kleinleed
en troostte haar met geduld.
Hij verschoonde Lieblings kleed,
wiens vreugde weer was vervuld.

De landheer werd geenszins boos,
maar vergaf haar slordigheid.
“Ach, er is maar weinig loos”,
zei hij vol barmhartigheid.

Toen verbeeldde Agape
dat ze waarlijk een mus was.
Hij wist dat op te rapen
en zag het niet als poespas,

maar speelde als vogel mee.
Ze fladderden in 't veld,
al keken anderen mee.
“Geen oordeel dat ons velt!”.

Ze klampte zich stevig vast
aan haar vaders grote hand.
Hij ervoer dat niet als last,
maar als een gewilde band.

Daarna zwaaide hij gedag,
zonder afscheid te pressen.
De ervaren landheer zag
haar strijd als levenslessen.

Bij haar deed hij niet te groot.
Was hij een mens die faalde.
Waarvan hij haar op zijn schoot
ook zonder schroom verhaalde.

Zo gingen dagen voorbij,
beschermd door warme rede.
Ze voelde zich blij en vrij,
in kalmte en in vrede.
